\documentclass[../../../include/open-logic-section]{subfiles}

\begin{document}

\olfileid{sth}{choice}{banach}
\olsection{બાનાખ--ટાર્સ્કીનો વિરોધાભાસ}

બૂલોઝે જેમ પ્રતિસ્થાપનને વાજબી
ઠરાવવાનો પ્રયાસ કર્યો, તેમ આપણે
પણ \emph{બાહ્ય} વિચારણાઓને અપીલ
કરીને પસંદગીને વાજબી ઠરાવવાનો
પ્રયત્ન કરી શકીએ
(\olref[replacement][extrinsic]{sec} જુઓ).
છેવટે પસંદગી સ્વીકારવાથી ઘણાં
ઇચ્છનીય પરિણામો મળે છે: બધા ગણાંકોની
તુલના કરી શકીએ; દરેક ગણને સુક્રમિત
કરી શકીએ; ગણાંકોને ક્રમવાચક
સંખ્યાઓના એક પ્રકાર તરીકે લઈ શકીએ;
વગેરે.

પરંતુ ક્યારેક એવો દાવો થાય છે કે
પસંદગીનાં \emph{અનિચ્છનીય} પરિણામો
છે. મોટાભાગે આ દાવો
\cite{BanachTarski1924}ના એક
પરિણામને કારણે થાય છે.

\begin{thm}[બાનાખ--ટાર્સ્કીનો વિરોધાભાસ ($\ZFC$માં)]
કોઈ પણ ઘન ગોળાને સાન્ત સંખ્યામાં
ટુકડાઓમાં વિભાજિત કરી શકાય છે.
તે ટુકડાઓ ફરી ગોઠવીને (પરિભ્રમણ
અને સ્થાનાંતરણથી) મૂળ ગોળાની
બે નકલો બનાવી શકાય છે.
\end{thm}
\noindent
પહેલી નજરે આ આશ્ચર્યજનક લાગે છે.
દેખીતી રીતે બે ગોળાઓનું ઘનફળ
મૂળ ગોળાના ઘનફળનું \emph{બમણું} છે.
પરંતુ અચળ આકારવાળી ગતિઓ---પરિભ્રમણ
અને સ્થાનાંતરણ---ઘનફળ બદલતી નથી.
એટલે બાનાખ--ટાર્સ્કીથી જાણે
નવું દ્રવ્ય જાદુથી અસ્તિત્વમાં
લાવી શકીએ એમ લાગે છે.

હજી વધુ આશ્ચર્ય થાય છે.\footnote{જુઓ
\citet[Theorem 3.12]{Wagon2016}.}
સમાન દલીલ બતાવે છે કે વટાણાને
સાન્ત સંખ્યામાં ટુકડાઓમાં કાપીને,
પછી તેમને (પરિભ્રમણ અને સ્થાનાંતરણથી)
ફરી ગોઠવી બિગ બેનના આકાર અને
કદની વસ્તુ બનાવી શકાય છે.

પરંતુ માત્ર $\ZF$માં આમાંથી
કંઈ સાબિત થતું નથી.\footnote{જોકે
પસંદગી કરતાં કડક રીતે નબળા
સિદ્ધાંતોમાંથી પણ બાનાખ--ટાર્સ્કી
સાબિત થાય છે; જુઓ
\citet[303]{Wagon2016}.}
તેથી નિર્ણય લેવો પડે: પસંદગી નકારવી,
અથવા ``વિરોધાભાસ'' સાથે રહેતાં
શીખવું.

અમે સૂચવીશું કે ``વિરોધાભાસ''
સાથે રહેતાં શીખવું જોઈએ. ખરેખર
અમને એ બહુ વિરોધાભાસી પણ લાગતો
નથી. ખાસ કરીને નીચેના કોઈ પણ
પરિણામ કરતાં તે વધુ કે ઓછો
વિરોધાભાસી શા માટે હોય તે સમજાતું
નથી:\footnote{\citet[276--7]{Potter2004},
\citet[16]{Weston2003} અને
\citet[31, 308--9]{Wagon2016}
અન્ય ઉદાહરણો વડે સમાન મુદ્દા કરે છે.}
\begin{enumerate}
	\item $(0,1)$ અંતરાલમાં $\Real$ જેટલાં જ બિંદુઓ છે.
	\\\emph{પુરાવો}: $\tan(\pi(r-\nicefrac{1}{2}))$ જુઓ.
	\item રેખામાં ચોરસ જેટલાં જ બિંદુઓ છે.
	\\\olref[his][set][pathology]{sec} અને
	\olref[his][set][cantorplane]{sec} જુઓ.
	\item અવકાશભર વક્રો છે.
	\\\olref[his][set][pathology]{sec} અને
	\olref[his][set][hilbertcurve]{sec} જુઓ.
\end{enumerate}
આ ત્રણ પરિણામોમાંથી કોઈ માટે
પસંદગી જરૂરી નથી. આજે આપણે
તેમને ગણિતનાં આશ્ચર્યજનક અને
સુંદર પરિણામો માનીએ છીએ.
કદાચ બાનાખ--ટાર્સ્કીના વિરોધાભાસ
પ્રત્યે પણ એ જ વલણ રાખવું જોઈએ.
\sourcecorrection{OLMD-150}{મૂળમાં
$\tan(\pi(r-\nicefrac{1}{2})))$
સૂત્રના અંતે એક વધારાનું બંધ
કૌંસ છે. અહીં તેને દૂર કરીને
$(0,1)$થી $\Real$ સુધીના
વિધેયનું યોગ્ય સૂત્ર લખ્યું છે.}

અલબત્ત, અહીં એક તકનીકી બાબત
ધ્યાનમાં લેવી પડે; પરંતુ શાંતિથી
વિચારીએ તો પૂરતું છે. અચળ આકારવાળી
ગતિઓ ઘનફળ જાળવે છે. તેથી ગોળાના
જે પાંચ\footnote{આપણે વિરોધાભાસને
``સાન્ત સંખ્યાના ટુકડાઓ'' કહીને
રજૂ કર્યો હતો. ખરેખર
\citet{Robinson1947}એ બતાવ્યું
કે વિભાજન \emph{પાંચ} ટુકડાઓથી
શક્ય છે (પણ તેનાથી ઓછાથી નહીં).
પુરાવા માટે \citet[pp.~66--7]{Wagon2016}
જુઓ.} ટુકડાઓ બને છે, તે બધા
\emph{માપનીય} હોઈ શકે નહીં.
અંદાજે કહીએ તો દરેક ટુકડાને અલગ
ઘનફળ આપવાનો અર્થ જ નથી. તેમને
બિંદુઓના ચિત્રમાં ન ઉતારી શકાય
એવા ``અનંત વિખેરાયેલા'' ગણો માનો.
આવા ગણોની કલ્પના ``વિચિત્ર'' લાગી
શકે. પરંતુ અગાઉ ઉપલબ્ધ ગણોના
\emph{બધા શક્ય} સંગ્રહો રચો એવો
\stagesacc{}માં રહેલો આદેશ તેમનું
અસ્તિત્વ સૂચવે છે.

આમાંથી કંઈ પણ મનાવે નહીં, તો
બાનાખ--ટાર્સ્કીનો વિરોધાભાસ
સ્વીકારવાના પક્ષમાં અંતિમ
(બાહ્ય) દલીલ છે. તેનાથી તરત જ
સર્વશ્રેષ્ઠ ગણિતીય રમૂજ મળે છે:
\begin{enumerate}
	\item[] \emph{પ્રશ્ન}. ``Banach--Tarski''નો અનાગ્રામ શું છે?
	\item[] \emph{જવાબ}. ``Banach--Tarski Banach--Tarski''.
\end{enumerate}

\end{document}
